% Fragmento público generado desde propietarios íntegros.
% Residencia: 52.6.3.
% Fuente: ../incoming/radion_monografia_20260827/manuscrito/sections/07_incertidumbre_oscilador.tex:112-155
\subsubsection{Oscilador inductivo--capacitivo asociado al par angular}

\paragraph{Construcción independiente de los 2 depósitos conjugados}

Dados \(\omega>0\) y una impedancia de referencia \(Z_*>0\), definimos
\begin{equation}
L_\eta=\frac{Z_*}{\omega}e^{-\eta},
\qquad
C_\eta=\frac{1}{Z_*\omega}e^\eta.
\label{eq:LC-def}
\end{equation}
Entonces
\begin{equation}
L_\eta C_\eta=\omega^{-2},
\qquad
\sqrt{\frac{L_\eta}{C_\eta}}=Z_*e^{-\eta}.
\label{eq:LC-invariants}
\end{equation}
La frecuencia permanece fija; la impedancia conserva la orientación de la
forma.

El Hamiltoniano es
\begin{equation}
H_{LC}=\frac{q^2}{2C_\eta}+\frac{\Phi^2}{2L_\eta}.
\label{eq:HLC}
\end{equation}
Con las variables normalizadas
\begin{equation}
Q=\sqrt{Z_*}\,q,\qquad
P=\frac{\Phi}{\sqrt{Z_*}},
\label{eq:LC-map}
\end{equation}
se obtiene
\begin{equation}
H_{LC}=\frac\omega2
\left(e^{-\eta}Q^2+e^\eta P^2\right).
\label{eq:HLC-QP}
\end{equation}
Sobre \eqref{eq:elipse-param},
\begin{equation}
H_{LC}=\hret\omega.
\label{eq:LC-level}
\end{equation}

% Fuente: ../incoming/radion_monografia_20260827/manuscrito/sections/07_incertidumbre_oscilador.tex:198-239
\paragraph{Oscilador LC como unidad dinámica elemental}

El oscilador armónico atraviesa casi toda la física porque linealiza el
entorno de un equilibrio, descompone campos en modos, organiza cuantos y
permite leer propagación como superposición de fases. En el radión, esa
universalidad recibe una genealogía concreta: el modo no empieza en una
ecuación diferencial externa, sino en una forma positiva con acción de vuelta
fija. La dinámica convencional aparece al elegir un reloj \(\omega\) y un
transductor \(Z_*\).

La frecuencia no altera la elipse de acción: cambia la energía del nivel
mediante \(E=\hret\omega\). Esta separación es crucial. La geometría fija la
acción; el reloj convierte acción en energía.

\subsubsection{Reversión temporal y propagaciones conjugadas}

Sea \(H=H^*\) y sea \(J_E\) una estructura compleja con \(J_E^2=-I\).
Supongamos un involutivo \(C\) tal que
\[
CJ_EC^{-1}=-J_E,\qquad CHC^{-1}=H.
\]
La evolución
\begin{equation}
U(t)=\exp\left(-\frac{J_EHt}{\hret}\right)
\label{eq:evolution-J}
\end{equation}
satisface exactamente
\begin{equation}
CU(t)C^{-1}=U(-t).
\label{eq:time-reversal}
\end{equation}
Ésta es la forma operatoria mínima de la bidireccionalidad: la conjugación
intercambia las prolongaciones temporalmente opuestas.

La teoría del absorbedor de Wheeler--Feynman y la lectura de antipartículas
como trayectorias temporalmente conjugadas ofrecen una realización histórica
posterior. La identidad \eqref{eq:time-reversal} está demostrada; convertirla
en una teoría completa de propagadores avanzados y retardados exige declarar
acción, condiciones de frontera, causalidad e interacciones del absorbedor.
La monografía conserva la intuición física sin convertir una correspondencia
en una igualdad de teorías.

